% Imporditud: tools/import_eksperimendid.py
% Allikas: Eesti füüsikaolümpiaadi eksperimendi ülesannete
%   kogu (Rael Kalda, 2023), ülesanne Ü52, lahendus L52.
\setAuthor{EFO žürii}
\setRound{lõppvoor}
\setYear{2008}
\setNumber{P E1}
\setDifficulty{2}
\setTopic{dünaamika}
\setVahendid{puidust klots, kumminiit, paberileht, mõõtejoonlaud}
\setPunktid{12}
\setAllikas{Eksperimendi kogu Ü52, lahendus lk 56}
\setLahendusseis{toores}

\prob{Hõõrdetegur}
Määrake liugehõõrdetegur klotsi ja paberi vahel.
\textit{Märkus:} Liugehõõrdetegur $\mu$ = hõõrdejõud / rõhumisjõud.

\hint

\solu
Riputame klotsi, mõõdame niidi pikenemise $L_1$ paneme klotsi alusele, tirime niidiga, mõõdame niidi pikenemise $L_2$. Et elastsusjõud on võrdeline niidi pikenemisega, siis $\mu$ = $F_2$/$F_1$ = $L_2$/$L_1$.
\probend
